\scanpage{128}
% archive foliation: 127
$X$, $P_X$, avec une \deleted{\unsure{immersion}}\inserted{connexion} intégrable, et une immersion \inserted{de $X$-espaces} (respectant les connexions)
\diagram{\begin{tikzcd}
\widetilde X\arrow[r,phantom,"\subset"]\arrow[dr,dash]&P_X\arrow[d,dash]\\
&X
\end{tikzcd}}
Comme $\widetilde X$ est muni d'une section marquée analytique, $P_X$ a une section marquée analytique.

Donc

\uline{Théorème 2} Pour $X$ courbe analytique \deleted{\unsure{décrite}}\inserted{non sing.}, on peut trouver, de façon unique :
(à iso unique près),
\[
\left\{\begin{aligned}
\text{a)}\quad&\text{Un fibré analytique }P_X\text{ sur }X\text{ en droites projectives}\\
\text{b)}\quad&\text{Une connexion (intégrable) sur }P_X\ (\text{rel : }X/\mathbb C)\\
\text{c)}\quad&\text{Une section }\sigma\text{ de }P_X\text{ sur }X\\
\text{d)}\quad&\text{Un ouvert }\widetilde X\text{ de }P_X\text{ contenant }\sigma(X)\\
\text{e)}\quad&\text{Un hom. de }X\text{-espaces analytiques }\widetilde X\xrightarrow{p}X_X=X\times X
\end{aligned}\right.
\]
\margin{Car b) c) impliquent déjà \unsure{un germe $\widetilde p$ d'hom. anal. prenant des iso transversaux} le long de $\sigma$ \unsure{supposée} (dire que $\sigma$ est nulle part stationnaire signifie que $p$ est un iso le long de $\sigma$) et on fait \unsure{alors du} prolongement analytique de ce germe ($\widetilde p\ldots$) cf \unsure{explication} plus bas.}
tels que
\[
\begin{aligned}
1^\circ)\quad&p\circ\sigma=\Delta_X\quad(\text{section diagonale})\\
2^\circ)\quad&\forall x\in X,\quad p_x:\widetilde X_x\longrightarrow(X_X)_x=X
\end{aligned}
\]
fait de $\widetilde X_x$ un rev. universel de $\inserted{\text{la comp. connexe de }x\text{ dans }}X$ (i.e. c'est un rev. de $X$, connexe et simplement connexe)

$3^\circ)$ $p$ est compatible avec les connexions

$4^\circ)$ $\forall x\in X$, $\widetilde X_x\subset(P_X)_x$ est un domaine homogène de $(P_X)_x$.

\margin{Les données \inserted{d/ et} e) satisfaisant $1^\circ\ 2^\circ\ 3^\circ$ sont déjà \uline{canoniquement} \unsure{définies par} a) b) c) \unsure{avec des}\ \uline{\unsure{restrictions}} dessus.}
